\documentclass[11pt]{article}

\usepackage[margin=1in]{geometry}
\usepackage{hyperref}
\usepackage{graphicx}
\usepackage{listings}
\usepackage{xcolor}
\usepackage{enumitem}
\usepackage{booktabs}

\hypersetup{
  pdftitle={multimodal-query-demo: an Arrow-native, tokio-based prototype},
  pdfauthor={Riaan de Beer},
  pdfsubject={Sparse multi-rate multimodal temporal event streams},
  pdfkeywords={Arrow, temporal, multimodal, query, Rust, tokio, async}
}

\lstset{
  basicstyle=\ttfamily\small,
  breaklines=true,
  columns=fullflexible,
  frame=single,
  framesep=4pt,
  xleftmargin=6pt,
  xrightmargin=6pt,
  showstringspaces=false,
}

\title{multimodal-query-demo:\\an Arrow-native, tokio-based prototype for\\sparse multi-rate multimodal temporal event streams}
\author{Riaan de Beer\\\texttt{predictiverendezvous@proton.me}}
\date{\today}

\begin{document}
\maketitle

\begin{abstract}
This design note accompanies a small Rust prototype for querying sparse multi-rate
multimodal temporal event streams. The prototype provides point-in-time and half-open
range scans over Apache Arrow \texttt{RecordBatch}es, two domain operators
(\emph{proximity\_pairs} and \emph{speed\_over\_ground}), and an async ingest pipeline
with snapshot-isolated reads via \texttt{arc\_swap::ArcSwap}. The artifact is
deliberately bounded: in-memory, single-process, single-runtime, synthetic workloads
only. A numbered non-claims charter is pinned both in source and in this paper, verified
by a test that fails CI if the two drift. The goal is a finished, technically honest
slice that demonstrates async-Rust judgment, Arrow-native data handling, and the kind of
query semantics that sparse multimodal logs demand, not a production system.
\end{abstract}

\section{Motivation}

Robotics, autonomous systems, and scientific instrumentation produce streams of events
whose modalities --- poses, scalars, annotations, image references, point clouds ---
update at independent cadences. A 200~Hz IMU, a 30~Hz camera, and a 1~Hz annotation
stream do not share a clock, and forcing them into a single row-aligned table introduces
either null-padding (wasted storage, confusing semantics) or resampling (information
loss before a query has even been posed). The temporal-database and time-series
literature has long preferred columnar, per-stream layouts with point-in-time
resolution; the prototype here takes that view and commits to it.

Two query shapes matter disproportionately in this setting. The first is
\emph{latest-at}: ``for each entity, what is its most recent pose as of time $t$?''
This is the query that lets a visualiser or a derived operator see a consistent frozen
world. The second is the \emph{range scan} with kind and entity pruning, which lets
derived operators walk a sparse slice of history without scanning the full log. Both
queries are cheap when the store is partitioned well and expensive when it is not; both
are easy to specify informally and surprisingly subtle to specify precisely --- ties at
equal timestamps, half-open vs.\ closed bounds, and snapshot isolation under concurrent
writes all need explicit contracts.

This prototype picks one concrete implementation path for those two queries, wraps it
in an async ingest pipeline, layers two representative domain operators on top
(\emph{proximity\_pairs} and \emph{speed\_over\_ground}), and measures the result
honestly. It is not trying to be a database. It is trying to show that the author can
make scope-sane design calls, write async Rust that earns its weight, and ship a thing
that builds, tests, and reproduces end-to-end.

\section{Scope and non-goals}
\label{sec:scope}

In scope: a four-crate Rust workspace (\texttt{mqd-core}, \texttt{mqd-engine},
\texttt{mqd-cli}, \texttt{mqd-bench}); five typed component kinds;
\texttt{latest\_at}\,/\,\texttt{range\_scan} with a query-plan surface and
\texttt{-{}-explain}; two domain operators; a deterministic seeded synthetic generator;
criterion benches plus a single-file \texttt{mqd bench} JSON report; figure rendering;
and this paper built by \texttt{latexmk}.

Out of scope, by design: durability (no WAL, no IPC files, no restart), distribution
(single process, single memory space), SQL or any third-party query engine, any network
surface, image or mesh byte-level decoding (we store references only), multi-timeline
joins beyond per-entity point-in-time resolution, and any claim of compatibility or
parity with any specific production system. The numbered non-claims charter in
Section~\ref{sec:nonclaims} fixes this boundary and is enforced by a CI test.

\section{Architecture}

A \texttt{tokio::sync::mpsc} channel moves \texttt{Event}s from one or more producer
tasks to a single flusher task, which assembles per-kind Arrow \texttt{RecordBatch}es
and appends them to a store partitioned by
\texttt{(ComponentKind, bucket\_start\_ns)}. The store keeps an
\texttt{arc\_swap::ArcSwap} pointer to the current \texttt{StoreSnapshot}. Readers
clone the pointer (lock-free) and hold an immutable view for the duration of their
query. A \texttt{tokio\_util::sync::CancellationToken} is threaded through the ingest
pipeline so that a final \texttt{flush\_all} fires exactly once, whether the CLI exits
cleanly or is cancelled mid-stream.

\begin{figure}[h]
\centering
\begin{tabular}{c}
\fbox{\parbox{0.9\textwidth}{\centering
\texttt{producers (N tasks)} $\longrightarrow$ \texttt{mpsc<Event>}
$\longrightarrow$ \texttt{flusher (1 task)} \\[2pt]
$\longrightarrow$ per-kind \texttt{KindBatcher} $\longrightarrow$
\texttt{Store::append\_batch} \\[2pt]
\texttt{ArcSwap<StoreSnapshot>} $\longleftarrow$ \texttt{swap} \\[2pt]
\texttt{queries / operators} $\longleftarrow$ \texttt{snapshot.load\_full()}
}}
\end{tabular}
\caption{Pipeline roles. One writer, many readers, snapshot isolated. The channel
bounds producer--flusher backpressure; the atomic swap is the only reader--writer
synchronisation point.}
\end{figure}

\section{Data model}

The store understands five \texttt{ComponentKind}s: \texttt{Transform3D},
\texttt{Points3D}, \texttt{Scalar}, \texttt{ImageRef}, and \texttt{Annotation}. Every
row carries the same two non-null prefix columns: \texttt{t\_ns:Int64} (event time, UTC
nanoseconds) and \texttt{entity\_path:Utf8} (a slash-delimited hierarchical identifier
such as \texttt{/agent/0}). The remaining columns depend on kind. A compact reference:

\begin{center}
\small
\begin{tabular}{@{}l l l@{}}
\toprule
\textbf{Kind} & \textbf{Columns} & \textbf{Notes} \\
\midrule
\texttt{Transform3D} & \texttt{tx,ty,tz,qx,qy,qz,qw:Float64} & translation + unit quaternion \\
\texttt{Points3D} & \texttt{points:List<Struct<x,y,z:Float64>{}>} & variable-length point cloud \\
\texttt{Scalar} & \texttt{value:Float64, unit:Utf8?} & \texttt{unit} nullable \\
\texttt{ImageRef} & \texttt{uri:Utf8, width:Int32, height:Int32, format:Utf8} & reference only, no pixels \\
\texttt{Annotation} & \texttt{label:Utf8} & string label \\
\bottomrule
\end{tabular}
\end{center}

Appendix~\ref{app:schemas} reproduces the full \texttt{arrow::datatypes::Schema}
definitions verbatim. The partition key is \texttt{(ComponentKind, bucket\_start\_ns)}
with a default bucket width of 1~second (\texttt{DEFAULT\_BUCKET\_NS = 10\textsuperscript{9}}).
Buckets are open-ended to the right: a batch whose time range spans a bucket boundary
is split into the minimum number of sub-batches needed to place each row in its correct
bucket. This keeps \texttt{range\_scan} pruning by bucket cheap and exact.

\section{Query semantics}

\begin{itemize}
  \item \textbf{latest\_at}\,$(t, \mathrm{kind}, \mathrm{entity}?)$: for each matching
    entity, the single row with the largest $t_{ns} \le t$. Ties at equal $t_{ns}$ are
    broken by the \emph{later-inserted} row, where insertion order is the lexicographic
    tuple \texttt{(partition\_ix, batch\_ix, row\_ix)}. The tie-break rule is
    verified by \texttt{latest\_at\_ties\_are\_deterministic}.
  \item \textbf{range\_scan}\,$(s, e, \mathrm{kind}, \mathrm{entity}?)$: every row
    with $s \le t_{ns} < e$ (half-open, matching Rust's \texttt{Range}). Inverted
    bounds ($s > e$) error.
  \item Both queries return a \texttt{QueryPlan} alongside the result batch:
    \texttt{candidate\_partitions}, \texttt{pruned\_by\_kind}, \texttt{pruned\_by\_time},
    \texttt{scanned\_partitions}, \texttt{scanned\_rows}, \texttt{returned\_rows},
    and \texttt{latency\_us}. Under \texttt{-{}-explain} the plan is serialised to the
    reporter. Example (terminal reporter):

\begin{lstlisting}
[mqd] query latest-at
[config]  entity=/agent/7 at_ns=2500000000 kind=Transform3D
[store]   snapshot_rows=80000 partitions=125
[plan]    candidate=125 pruned_by_kind=100 pruned_by_time=22 scanned=3
[exec]    scanned_rows=1536 returned_rows=1
[result]  latency=41us
[done]    ok
\end{lstlisting}
\end{itemize}

Snapshot isolation: a reader's \texttt{Arc<StoreSnapshot>} is frozen for the duration
of the query regardless of concurrent writes. The
\texttt{tests/concurrent\_snapshots.rs} integration test spawns a background writer and
asserts that readers never observe a decreasing row count, that every snapshot is
internally consistent, and that the final snapshot after cancellation contains every
event the producer emitted.

\section{Domain operators}

\paragraph{proximity\_pairs.} At a point in time $t$, resolve each entity's latest
\texttt{Transform3D} pose (internally a \texttt{latest\_at} call --- one source of
truth for point-in-time semantics), then emit every unordered pair whose pose
separation is $\le r$. Complexity is $O(N^2)$ in the number of distinct entities at
$t$; no spatial index is used. This is the intended shape and is confirmed, not
refuted, by the benchmark in Figure~\ref{fig:proximity}. Output schema:
\texttt{at\_ns:Int64, entity\_a:Utf8, entity\_b:Utf8, distance\_m:Float64}, sorted
lexicographically on \texttt{(entity\_a, entity\_b)}.

\paragraph{speed\_over\_ground.} For one entity within an optional half-open window,
select every \texttt{Transform3D} row (internally a \texttt{range\_scan} call), stable-
sort by $t_{ns}$, and emit one row per consecutive pair with \texttt{dt\_ns},
\texttt{distance\_m}, and $\mathrm{speed\_mps} = \mathrm{distance\_m}/\Delta t$. Pairs
with $\Delta t \le 0$ are skipped (duplicate-timestamp noise rather than a semantic
event). Output schema: \texttt{entity\_path:Utf8, t\_start\_ns:Int64, t\_end\_ns:Int64,
dt\_ns:Int64, distance\_m:Float64, speed\_mps:Float64}.

Correctness fixtures: \texttt{speed\_constant\_one\_meter\_per\_second} emits a
10-tick trajectory of 1~m/s due east and asserts every consecutive-pair speed is
within $10^{-9}$ of 1.0. \texttt{proximity\_symmetric\_pair\_exactly\_on\_radius}
verifies boundary behaviour and lexicographic ordering.

\section{Async design}

Three concurrent roles share one snapshot:

\begin{itemize}
  \item \textbf{Producers} (N tasks, one per input source or synthetic generator
    partition) push \texttt{Event}s into a bounded \texttt{mpsc} channel.
  \item \textbf{Flusher} (1 task) \texttt{recv()}s into per-kind
    \texttt{KindBatcher}s, emits \texttt{RecordBatch}es when either
    \texttt{max\_rows} or \texttt{max\_delay} is hit, and calls
    \texttt{Store::append\_batch}, which constructs a new \texttt{StoreSnapshot} and
    atomically replaces the \texttt{ArcSwap} pointer.
  \item \textbf{Readers} (any task) call \texttt{store.snapshot()}, which is
    \texttt{O(1)} and lock-free, and then operate on the frozen snapshot for the
    lifetime of their query.
\end{itemize}

A \texttt{tokio\_util::sync::CancellationToken} is shared with the flusher. Its
selector uses \texttt{tokio::select!\ biased} to prefer cancellation over new work,
then drains any events remaining in the channel, then calls \texttt{flush\_all} exactly
once before returning. This is the property tested by
\texttt{tests/concurrent\_snapshots.rs}: after a cancellation, the final snapshot
contains every event the producer emitted, no more, no less.

What this design \emph{does} prove: that a single-writer, many-reader, snapshot-isolated
pattern is natural to express in async Rust given the right primitives (\texttt{mpsc},
\texttt{CancellationToken}, \texttt{ArcSwap}), and that writing an integration test
against a multi-threaded tokio runtime is straightforward. What it does \emph{not}
prove: correctness under adversarial concurrent loads, memory-ordering edge cases beyond
what \texttt{ArcSwap} already guarantees, or performance under realistic fan-in.

\section{Implementation notes}

\paragraph{Partition key.} \texttt{(ComponentKind, bucket\_start\_ns)}. Kind is a
hard boundary (every query specifies a kind, so we prune a huge fraction of partitions
at the top); bucket is a soft-sharded time axis. Default bucket width is $10^9$~ns
(1~s); this is exported as \texttt{DEFAULT\_BUCKET\_NS} and fed into
\texttt{Store::new}. A single input \texttt{RecordBatch} whose \texttt{t\_ns} range
spans a bucket boundary is split at the boundary into the minimum number of sub-batches
needed, so every resulting partition entry has a $[\min t_{ns}, \max t_{ns}]$ envelope
fully contained in a single bucket. This makes bucket-level time pruning exact and
lets \texttt{range\_scan} skip whole partitions by arithmetic alone.

\paragraph{Batch sizing.} The CLI exposes \texttt{-{}-batch-rows} (default~4096). This
trades per-batch fixed overhead (schema resolution, appending to partition list,
snapshot rebuild) against memory footprint per batch. The default is in the range where
Arrow's columnar overhead amortises well; no automated tuning is claimed.

\paragraph{RNG determinism.} \texttt{rand\_pcg::Pcg64} is seeded from \texttt{-{}-seed}.
The generator emits events in strict tick order and, within a tick, entities in sorted
\texttt{EntityPath} order --- a \texttt{BTreeMap} iteration, never a
\texttt{HashMap}. Two fingerprint tests pin the output byte-for-byte: the JSONL
produced by \texttt{mqd generate --seed 42 --agents 4 --ticks 100}, and the
\texttt{proximity\_pairs} result at \texttt{at=250\_000\_000, radius=2.5} on the
corresponding store. These are SHA-256 hashes of the raw columnar data --- not of
Arrow IPC bytes --- so the test is robust to Arrow version encoding changes while
still pinning every user-visible value.

\paragraph{Tie-break contract.} Equal \texttt{t\_ns} values route to the same partition
(same bucket), so within a tie set the lexicographic tuple
\texttt{(partition\_ix, batch\_ix, row\_ix)} is a total order. \texttt{latest\_at}
scans candidate partitions in this order and keeps the last match when timestamps
tie --- giving the ``later-inserted row wins'' guarantee for free, without a dedicated
sequence column.

\paragraph{Error surface.} \texttt{mqd-core} exposes a small \texttt{Error} enum via
\texttt{thiserror}; user-facing \texttt{anyhow::Result} is used at binary boundaries
only. Engine-internal calls never use \texttt{anyhow}. Schema mismatches, inverted
range bounds, and malformed entity paths are all typed errors, not panics.

\section{Prototype evaluation}

The figures below are rendered from a single \texttt{bench.json} report produced by
\texttt{mqd bench}. Numbers are single-machine wall-clock latencies on one developer
workstation; see Section~\emph{How to read} for the caveats they carry.

\begin{figure}[h]
  \centering
  \includegraphics[width=0.7\textwidth]{figures/ingest_latency.png}
  \caption{Ingest (batch + store append) latency versus total event count. Solid line
  p50, dashed line p95.}
\end{figure}

\begin{figure}[h]
  \centering
  \includegraphics[width=0.7\textwidth]{figures/latest_at_latency.png}
  \caption{\texttt{latest\_at} latency across agent counts at fixed tick budget. Solid
  p50, dashed p95.}
\end{figure}

\begin{figure}[h]
  \centering
  \includegraphics[width=0.7\textwidth]{figures/range_scan_latency.png}
  \caption{\texttt{range\_scan} latency for a 25\% window on one entity. Solid p50,
  dashed p95.}
\end{figure}

\begin{figure}[h]
  \label{fig:proximity}
  \centering
  \includegraphics[width=0.7\textwidth]{figures/proximity_latency.png}
  \caption{\texttt{proximity\_pairs} latency on log-log axes. Solid p50, dashed p95,
  grey dashed $O(N^2)$ reference anchored at the smallest measured point. The
  quadratic shape is \emph{confirmed} by approximate slope-2 on log-log, not refuted.}
\end{figure}

\begin{figure}[h]
  \centering
  \includegraphics[width=0.7\textwidth]{figures/speed_latency.png}
  \caption{\texttt{speed\_over\_ground} latency on log-log axes with an $O(N)$
  reference line. Solid p50, dashed p95.}
\end{figure}

\section{How to read the prototype outputs}

All latencies are wall-clock on a single machine under a multi-threaded tokio runtime.
Ingest numbers measure Arrow-shaped event production through the channel and snapshot
swap; JSONL parsing is excluded. Proximity scaling is quadratic by construction. These
numbers are prototype signal, not product benchmarks, and cannot be extrapolated to any
other workload or system. Variance between runs on the same workstation is not
characterised. The sample count (\texttt{-{}-samples}, default 51) is the smallest $n$
for which nearest-rank p95 and p99 land on distinct order statistics
(\texttt{sorted[48]} and \texttt{sorted[50]} respectively); the figures show p50 and
p95 only because at this $n$ p99 is equal to the maximum observation and carries no
more information than "worst in the run". The \texttt{-{}-samples} flag lets the
reader increase $n$ for a narrower percentile estimate.

\section{Limitations}

Synthetic data only. No spatial index (proximity is $O(N^2)$ by design). Concurrency
tested via unit and integration tests only; adversarial concurrent loads have not been
stress-tested. Memory has not been profiled. Semantic correctness is bounded by fixture
coverage. No durability: the store lives entirely in the process heap and is lost on
exit. No multi-timeline joins: every query works against a single event-time timeline,
and the \texttt{Timeline} type exists but is not exercised beyond that.

\section{Non-claims}
\label{sec:nonclaims}

\begin{enumerate}[label=\arabic*.]
  \item multimodal-query-demo is not a database. It is not durable, not transactional, and does not survive process restart.
  \item multimodal-query-demo does not claim compatibility or parity with any specific production system. It is an independent design exercise against public literature on sparse multimodal temporal logs.
  \item Point-in-time query semantics are validated only against the synthetic fixtures in the test suite. Correctness against arbitrary real-world workloads is not claimed.
  \item multimodal-query-demo does not implement SQL, a third-party query engine, a network protocol, or any query interface beyond the typed operators exposed on its CLI.
  \item Benchmark numbers are single-machine, synthetic-workload latencies on the author's hardware. They are not throughput claims, not SLA claims, and not comparisons against any other system.
  \item multimodal-query-demo does not claim cloud-native, distributed, or multi-node behavior. It is one process, one memory space, orchestrated by a single tokio runtime.
\end{enumerate}

\section{Reproducibility}

\begin{lstlisting}
cargo test --workspace
cargo run --release -p mqd-cli -- bench --out out/bench.json
cargo run --release -p mqd-cli -- figs --from out/bench.json \
    --out paper/figures/
./paper/build.sh
\end{lstlisting}

The deterministic-replay test pins SHA-256 fingerprints for (a) the JSONL generator
output at \texttt{seed=42, agents=4, ticks=100}, and (b) the
\texttt{proximity\_pairs} result on the corresponding store at
\texttt{at=250\_000\_000, radius=2.5}. The non-claims lock test parses this paper's
\texttt{Non-claims} enumerate block and asserts byte-for-byte equality with
\texttt{mqd\_core::NON\_CLAIMS}. Both tests fail CI on any drift.

\section{Future work}

Plausible, small extensions only --- this is a sketch of where a next iteration would
go, not a roadmap.

\begin{itemize}
  \item \textbf{Durability via Arrow IPC.} A \texttt{-{}-dump out.arrow} path that
    serialises every partition's batch list to an Arrow IPC file, and a loader that
    hydrates a \texttt{Store} from one. No write-ahead log; no crash consistency.
  \item \textbf{Spatial index for proximity.} A $k$-d tree or uniform grid rebuilt per
    call would drop \texttt{proximity\_pairs} from $O(N^2)$ to expected $O(N \log N)$
    for reasonable densities. The operator is already isolated behind a \texttt{Query}
    type; swapping the implementation is a local change.
  \item \textbf{Richer time-coalescing.} A \texttt{latest\_within} operator that
    returns the latest row within a backward window rather than the single nearest
    $t_{ns} \le t$, useful when callers want a ``freshness budget'' rather than a
    point.
  \item \textbf{Bounded-memory range scans.} Stream \texttt{RecordBatch}es out of
    \texttt{range\_scan} through an \texttt{impl Stream<Item = RecordBatch>} rather
    than materialising one batch, for scans large enough that the batch itself is the
    bottleneck.
\end{itemize}

\appendix

\section{CLI reference}
\label{app:cli}

\begin{lstlisting}
mqd generate   --seed <u64> --agents <u32> --ticks <u32> --out <path>
mqd ingest     --input <path> [--batch-rows <usize>]
mqd query latest-at --input <path> --at <i64> [--entity <path>]
                    [--kind <kind>] [--explain]
mqd query range     --input <path> --start <i64> --end <i64>
                    [--entity <path>] [--kind <kind>] [--explain]
mqd op proximity    --input <path> --at <i64> --radius <f64>
                    [--explain]
mqd op speed        --input <path> --entity <path>
                    [--start <i64>] [--end <i64>] [--explain]
mqd bench           [--out <path>] [--samples <usize>]
mqd figs            --from <path> --out <dir>
mqd non-claims
\end{lstlisting}

Global: \texttt{-{}-json} before any subcommand switches the reporter from the
terminal formatter to newline-delimited JSON (one event per line, suitable for
\texttt{jq}).

\section{Arrow schemas}
\label{app:schemas}

All schemas share the prefix
\texttt{t\_ns: Int64 (non-null), entity\_path: Utf8 (non-null)}.

\begin{lstlisting}
Transform3D:
  t_ns: Int64, entity_path: Utf8,
  tx: Float64, ty: Float64, tz: Float64,
  qx: Float64, qy: Float64, qz: Float64, qw: Float64

Points3D:
  t_ns: Int64, entity_path: Utf8,
  points: List<Struct<x: Float64, y: Float64, z: Float64>>

Scalar:
  t_ns: Int64, entity_path: Utf8,
  value: Float64, unit: Utf8 (nullable)

ImageRef:
  t_ns: Int64, entity_path: Utf8,
  uri: Utf8, width: Int32, height: Int32, format: Utf8

Annotation:
  t_ns: Int64, entity_path: Utf8,
  label: Utf8
\end{lstlisting}

All non-prefix fields are non-null unless marked otherwise. The canonical definitions
live in \texttt{crates/mqd-core/src/schema.rs}; changes to them flow through a single
\texttt{schema\_for(kind)} function consumed by the engine, the CLI, and the bench
harness.

\end{document}
